\section{Note on Sources and Delivery}

This homily addresses an adult parish assembly for the Resumed Sixth Sunday after the Epiphany, 15 November 2026, in the 1962 Roman Missal. Its spoken body contains \textbf{1,185 words}, excluding the title and this terminal apparatus. At an unhurried 115 words per minute it takes approximately \textbf{10 minutes 18 seconds}; at 110--120 words per minute, approximately 9 minutes 53 seconds to 10 minutes 46 seconds. These are arithmetic estimates, not a timed human delivery. A silent text review checked the argument, sentence clarity, paragraph endings and transitions; no audible rehearsal is claimed.

The homily follows the reviewed study's interpretation \textit{The word received in trial becomes a witness to others}, as compared in the concise study. Chrysostom's explanation of leaven working through contact and Aquinas's explanation of exemplary faith under trial supply its principal reasoning. The connections with the Collect, Secret, repeated psalm petition and Postcommunion, and the application to patient speech, are the homily's own synthesis from that interpretation. The decisive detail is the leaven hidden within the meal, read beside the Thessalonians' becoming examples while still afflicted. Difficulty therefore need not postpone faithful action or its benefit to another believer. The speech preserves the distinction between suffering and the Spirit's joy, and between present example and the resurrection still awaited.

The speech is authored preaching prose, not an official liturgical text. Its account of the prayers is commentary rather than a text for recitation. The historical sources retain their own rights; no modern translation or editorial expression is reproduced. The canonical research records identify the exact source witnesses and their limits.

\section{References}

\begin{itemize}
\item \textbf{Appointed texts and biblical contexts.} \textit{Missale Romanum}, Vatican typical edition, 1962, pp.~421--422, nos.~1759--1768: especially the Collect, no.~1760; Epistle, no.~1761; Alleluia, no.~1763; Gospel, no.~1764; Offertory, no.~1765; Secret, no.~1766; and Postcommunion, no.~1768. \url{https://media.churchmusicassociation.org/pdf/missale62.pdf}. Biblical context: Matthew 13:1--33; 1 Thessalonians 1:2--10; Psalm 129:1--2 in Vulgate numbering (Hebrew Psalm 130), read in the registered Douay--Rheims (Challoner), Project Gutenberg no.~1581. The Epistle's body corrects the erroneous citation printed on p.~421; the corresponding winter lesson is p.~49, no.~338.

\item \textbf{John Chrysostom.} \textit{Homilies on Matthew}, 46.2--3, on the preceding parables, mustard seed, leaven in contact with the meal and apostolic virtue; historical English in \textit{Nicene and Post-Nicene Fathers}, first series, vol.~10, ed. Philip Schaff, CCEL text. \textit{Homilies on First Thessalonians}, homily 1, on 1:2--7, especially the Spirit's joy within affliction; NPNF first series, vol.~13, CCEL text, with the Thessalonian translation attributed in its preface to James Tweed. No new collation of the Greek is claimed.

\item \textbf{Thomas Aquinas.} \textit{Super I Epistolam ad Thessalonicenses lectura}, chapter 1, lecture 1 [87862], especially on 1:3, 1:6--9 and 1:10: active faith, hope's endurance, imitation becoming example, resurrection and deliverance. Turin 1953 Latin edition, Corpus Thomisticum delivery revised by Enrique Alarc\'on, \url{https://www.corpusthomisticum.org/c1t.html}. Used through original attributed exposition; the restricted modern editorial text is not reproduced.
\end{itemize}
